% --- Archon physics formalization source begin ---
% archon:physics
% archon:covers PhyXMiniProblems/problem_phyx_mini_0288.lean
% archon:source-report reports/phyx_mini/problem_phyx_mini_0288.source.json

\chapter{Physics problem phyx\_mini\_0288}
\label{ch:PhyXMiniProblems_problem_phyx_mini_0288}

\paragraph{Problem source.}
\#\# Physical scenario

Figure shows the transverse velocity \$u\$ versus time \$t\$ of the point on a string at \$x = 0\$, as a wave passes through it. The scale on the vertical axis is set by \$u\_s = 4.0\$ m/s. The wave has the generic form \$y(x, t) = y\_m \textbackslash{}sin(kx - \textbackslash{}omega t + \textbackslash{}phi)\$.

\#\# Question

What then is \$\textbackslash{}phi\$?

\#\# Answer choices

- A: A: -0.78 rad

- B: B: -0.36 rad

- C: C: 0.64 rad

- D: D: -0.64 rad

\#\# Figure caption (auxiliary; use the image as primary evidence)

The image is a graph with a sine wave-like curve. The x-axis is labeled as "t", and the y-axis is labeled as "u (m/s)". The wave starts below the x-axis, rises through it, reaches a peak, and then descends back through the axis. The y-axis has markings at two points: "us" and "-us", indicating the peak positive and negative values of the curve. Grid lines are present to aid in reading the graph. The curve likely represents a periodic function of speed over time.

\#\# Dataset metadata

Wave Properties; Physical Model Grounding Reasoning, Spatial Relation Reasoning

\paragraph{Recorded answer/context.}
D: D: -0.64 rad

\paragraph{Figure/image path.}
/root/proposal\_for\_physic/science-mango/phyx\_mini\_run/phyx\_data/test\_image/288.png

\paragraph{Formalization target.}
create a compiling Lean file with sorry bodies at `PhyXMiniProblems/problem\_phyx\_mini\_0288.lean`. The Lean declarations must preserve the physical quantities, dimensions or dimensional roles, figure labels, governing-law hypotheses, and final relation expressed by this problem.
Use Mathlib/Physlib names found through LeanExplore where available. If a physics API is missing, introduce faithful local abstractions rather than scalar placeholder aliases.

\begin{theorem}[Physics formalization target]
\label{thm:physics:phyx_mini_0288:target}
\lean{PhyXMiniProblems.ProblemPhyXMini0288.problem_phyx_mini_0288}
\uses{def:physics:phyx-mini-0288:phyxminiproblems-problemphyxmini0288-sinusoidalstringwavesetup, def:physics:phyx-mini-0288:phyxminiproblems-problemphyxmini0288-matchestransversevelocitygraph, def:physics:phyx-mini-0288:phyxminiproblems-problemphyxmini0288-satisfiessinusoidaltravelingwave, def:physics:phyx-mini-0288:phyxminiproblems-problemphyxmini0288-recordeddatasetanswer, def:physics:phyx-mini-0288:phyxminiproblems-problemphyxmini0288-isnearestdisplayedphase}
The assigned autoformalize agent should translate this physics problem into Lean declarations in the covered file, with theorem and lemma proof bodies written as `by sorry`.
\end{theorem}
\begin{proof}
This is an autoformalization task, not a proof task. Produce statements that can later be proved by the physics prover without weakening the physical model.
\end{proof}

% --- Archon declaration topology begin ---
\section{Declaration topology}
\begin{definition}[velocity Dimension]
  \label{def:physics:phyx-mini-0288:phyxminiproblems-problemphyxmini0288-velocitydimension}
  \lean{PhyXMiniProblems.ProblemPhyXMini0288.velocityDimension}
  Velocity has physical dimension length divided by time
\end{definition}
\begin{proof}
  This is the declared model definition of velocity Dimension.
\end{proof}
\begin{definition}[acceleration Dimension]
  \label{def:physics:phyx-mini-0288:phyxminiproblems-problemphyxmini0288-accelerationdimension}
  \lean{PhyXMiniProblems.ProblemPhyXMini0288.accelerationDimension}
  Acceleration has physical dimension length divided by time squared
\end{definition}
\begin{proof}
  This is the declared model definition of acceleration Dimension.
\end{proof}
\begin{definition}[Length Magnitude]
  \label{def:physics:phyx-mini-0288:phyxminiproblems-problemphyxmini0288-lengthmagnitude}
  \lean{PhyXMiniProblems.ProblemPhyXMini0288.LengthMagnitude}
  A nonnegative displacement amplitude
\end{definition}
\begin{proof}
  This is the declared model definition of Length Magnitude.
\end{proof}
\begin{definition}[Signed Length Quantity]
  \label{def:physics:phyx-mini-0288:phyxminiproblems-problemphyxmini0288-signedlengthquantity}
  \lean{PhyXMiniProblems.ProblemPhyXMini0288.SignedLengthQuantity}
  A signed coordinate or transverse displacement on the one-dimensional string
\end{definition}
\begin{proof}
  This is the declared model definition of Signed Length Quantity.
\end{proof}
\begin{definition}[Wave Number Magnitude]
  \label{def:physics:phyx-mini-0288:phyxminiproblems-problemphyxmini0288-wavenumbermagnitude}
  \lean{PhyXMiniProblems.ProblemPhyXMini0288.WaveNumberMagnitude}
  A nonnegative wave-number magnitude, carrying inverse-length dimension
\end{definition}
\begin{proof}
  This is the declared model definition of Wave Number Magnitude.
\end{proof}
\begin{definition}[Angular Frequency Magnitude]
  \label{def:physics:phyx-mini-0288:phyxminiproblems-problemphyxmini0288-angularfrequencymagnitude}
  \lean{PhyXMiniProblems.ProblemPhyXMini0288.AngularFrequencyMagnitude}
  A nonnegative angular-frequency magnitude, carrying inverse-time dimension
\end{definition}
\begin{proof}
  This is the declared model definition of Angular Frequency Magnitude.
\end{proof}
\begin{definition}[Speed Magnitude]
  \label{def:physics:phyx-mini-0288:phyxminiproblems-problemphyxmini0288-speedmagnitude}
  \lean{PhyXMiniProblems.ProblemPhyXMini0288.SpeedMagnitude}
  \uses{def:physics:phyx-mini-0288:phyxminiproblems-problemphyxmini0288-velocitydimension}
  A nonnegative transverse-speed magnitude
\end{definition}
\begin{proof}
  This is the declared model definition of Speed Magnitude.
\end{proof}
\begin{definition}[Signed Velocity Quantity]
  \label{def:physics:phyx-mini-0288:phyxminiproblems-problemphyxmini0288-signedvelocityquantity}
  \lean{PhyXMiniProblems.ProblemPhyXMini0288.SignedVelocityQuantity}
  \uses{def:physics:phyx-mini-0288:phyxminiproblems-problemphyxmini0288-velocitydimension}
  A signed transverse-velocity component
\end{definition}
\begin{proof}
  This is the declared model definition of Signed Velocity Quantity.
\end{proof}
\begin{definition}[Signed Acceleration Quantity]
  \label{def:physics:phyx-mini-0288:phyxminiproblems-problemphyxmini0288-signedaccelerationquantity}
  \lean{PhyXMiniProblems.ProblemPhyXMini0288.SignedAccelerationQuantity}
  \uses{def:physics:phyx-mini-0288:phyxminiproblems-problemphyxmini0288-accelerationdimension}
  A signed transverse-acceleration component
\end{definition}
\begin{proof}
  This is the declared model definition of Signed Acceleration Quantity.
\end{proof}
\begin{definition}[length Magnitude Readout]
  \label{def:physics:phyx-mini-0288:phyxminiproblems-problemphyxmini0288-lengthmagnitudereadout}
  \lean{PhyXMiniProblems.ProblemPhyXMini0288.lengthMagnitudeReadout}
  \uses{def:physics:phyx-mini-0288:phyxminiproblems-problemphyxmini0288-lengthmagnitude}
  Read a nonnegative length in a selected length unit
\end{definition}
\begin{proof}
  This is the declared model definition of length Magnitude Readout.
\end{proof}
\begin{definition}[signed Length Readout]
  \label{def:physics:phyx-mini-0288:phyxminiproblems-problemphyxmini0288-signedlengthreadout}
  \lean{PhyXMiniProblems.ProblemPhyXMini0288.signedLengthReadout}
  \uses{def:physics:phyx-mini-0288:phyxminiproblems-problemphyxmini0288-signedlengthquantity}
  Read a signed length in a selected length unit
\end{definition}
\begin{proof}
  This is the declared model definition of signed Length Readout.
\end{proof}
\begin{definition}[wave Number Readout]
  \label{def:physics:phyx-mini-0288:phyxminiproblems-problemphyxmini0288-wavenumberreadout}
  \lean{PhyXMiniProblems.ProblemPhyXMini0288.waveNumberReadout}
  \uses{def:physics:phyx-mini-0288:phyxminiproblems-problemphyxmini0288-wavenumbermagnitude}
  Read wave number in inverse units of a selected length unit
\end{definition}
\begin{proof}
  This is the declared model definition of wave Number Readout.
\end{proof}
\begin{definition}[angular Frequency Readout]
  \label{def:physics:phyx-mini-0288:phyxminiproblems-problemphyxmini0288-angularfrequencyreadout}
  \lean{PhyXMiniProblems.ProblemPhyXMini0288.angularFrequencyReadout}
  \uses{def:physics:phyx-mini-0288:phyxminiproblems-problemphyxmini0288-angularfrequencymagnitude}
  Read angular frequency in inverse units of a selected time unit
\end{definition}
\begin{proof}
  This is the declared model definition of angular Frequency Readout.
\end{proof}
\begin{definition}[speed Readout]
  \label{def:physics:phyx-mini-0288:phyxminiproblems-problemphyxmini0288-speedreadout}
  \lean{PhyXMiniProblems.ProblemPhyXMini0288.speedReadout}
  \uses{def:physics:phyx-mini-0288:phyxminiproblems-problemphyxmini0288-speedmagnitude}
  Read a nonnegative speed in coherent selected units
\end{definition}
\begin{proof}
  This is the declared model definition of speed Readout.
\end{proof}
\begin{definition}[signed Velocity Readout]
  \label{def:physics:phyx-mini-0288:phyxminiproblems-problemphyxmini0288-signedvelocityreadout}
  \lean{PhyXMiniProblems.ProblemPhyXMini0288.signedVelocityReadout}
  \uses{def:physics:phyx-mini-0288:phyxminiproblems-problemphyxmini0288-signedvelocityquantity}
  Read a signed velocity in coherent selected units
\end{definition}
\begin{proof}
  This is the declared model definition of signed Velocity Readout.
\end{proof}
\begin{definition}[signed Acceleration Readout]
  \label{def:physics:phyx-mini-0288:phyxminiproblems-problemphyxmini0288-signedaccelerationreadout}
  \lean{PhyXMiniProblems.ProblemPhyXMini0288.signedAccelerationReadout}
  \uses{def:physics:phyx-mini-0288:phyxminiproblems-problemphyxmini0288-signedaccelerationquantity}
  Read a signed acceleration in coherent selected units
\end{definition}
\begin{proof}
  This is the declared model definition of signed Acceleration Readout.
\end{proof}
\begin{definition}[length In Meters]
  \label{def:physics:phyx-mini-0288:phyxminiproblems-problemphyxmini0288-lengthinmeters}
  \lean{PhyXMiniProblems.ProblemPhyXMini0288.lengthInMeters}
  \uses{def:physics:phyx-mini-0288:phyxminiproblems-problemphyxmini0288-signedlengthquantity, def:physics:phyx-mini-0288:phyxminiproblems-problemphyxmini0288-signedlengthreadout}
  Read a signed position or displacement in metres
\end{definition}
\begin{proof}
  This is the declared model definition of length In Meters.
\end{proof}
\begin{definition}[wave Number In Radians Per Meter]
  \label{def:physics:phyx-mini-0288:phyxminiproblems-problemphyxmini0288-wavenumberinradianspermeter}
  \lean{PhyXMiniProblems.ProblemPhyXMini0288.waveNumberInRadiansPerMeter}
  \uses{def:physics:phyx-mini-0288:phyxminiproblems-problemphyxmini0288-wavenumbermagnitude, def:physics:phyx-mini-0288:phyxminiproblems-problemphyxmini0288-wavenumberreadout}
  Read wave number in radians per metre
\end{definition}
\begin{proof}
  This is the declared model definition of wave Number In Radians Per Meter.
\end{proof}
\begin{definition}[angular Frequency In Radians Per Second]
  \label{def:physics:phyx-mini-0288:phyxminiproblems-problemphyxmini0288-angularfrequencyinradianspersecond}
  \lean{PhyXMiniProblems.ProblemPhyXMini0288.angularFrequencyInRadiansPerSecond}
  \uses{def:physics:phyx-mini-0288:phyxminiproblems-problemphyxmini0288-angularfrequencymagnitude, def:physics:phyx-mini-0288:phyxminiproblems-problemphyxmini0288-angularfrequencyreadout}
  Read angular frequency in radians per second
\end{definition}
\begin{proof}
  This is the declared model definition of angular Frequency In Radians Per Second.
\end{proof}
\begin{definition}[speed In Meters Per Second]
  \label{def:physics:phyx-mini-0288:phyxminiproblems-problemphyxmini0288-speedinmeterspersecond}
  \lean{PhyXMiniProblems.ProblemPhyXMini0288.speedInMetersPerSecond}
  \uses{def:physics:phyx-mini-0288:phyxminiproblems-problemphyxmini0288-speedmagnitude, def:physics:phyx-mini-0288:phyxminiproblems-problemphyxmini0288-speedreadout}
  Read a nonnegative speed in metres per second
\end{definition}
\begin{proof}
  This is the declared model definition of speed In Meters Per Second.
\end{proof}
\begin{definition}[velocity In Meters Per Second]
  \label{def:physics:phyx-mini-0288:phyxminiproblems-problemphyxmini0288-velocityinmeterspersecond}
  \lean{PhyXMiniProblems.ProblemPhyXMini0288.velocityInMetersPerSecond}
  \uses{def:physics:phyx-mini-0288:phyxminiproblems-problemphyxmini0288-signedvelocityquantity, def:physics:phyx-mini-0288:phyxminiproblems-problemphyxmini0288-signedvelocityreadout}
  Read a signed transverse velocity in metres per second
\end{definition}
\begin{proof}
  This is the declared model definition of velocity In Meters Per Second.
\end{proof}
\begin{definition}[acceleration In Meters Per Second Squared]
  \label{def:physics:phyx-mini-0288:phyxminiproblems-problemphyxmini0288-accelerationinmeterspersecondsquared}
  \lean{PhyXMiniProblems.ProblemPhyXMini0288.accelerationInMetersPerSecondSquared}
  \uses{def:physics:phyx-mini-0288:phyxminiproblems-problemphyxmini0288-signedaccelerationquantity, def:physics:phyx-mini-0288:phyxminiproblems-problemphyxmini0288-signedaccelerationreadout}
  Read a signed transverse acceleration in metres per second squared
\end{definition}
\begin{proof}
  This is the declared model definition of acceleration In Meters Per Second Squared.
\end{proof}
\begin{definition}[Graph Axis Symbol]
  \label{def:physics:phyx-mini-0288:phyxminiproblems-problemphyxmini0288-graphaxissymbol}
  \lean{PhyXMiniProblems.ProblemPhyXMini0288.GraphAxisSymbol}
  The literal symbols printed beside the graph axes
\end{definition}
\begin{proof}
  This is the declared model definition of Graph Axis Symbol.
\end{proof}
\begin{definition}[Velocity Graph Label]
  \label{def:physics:phyx-mini-0288:phyxminiproblems-problemphyxmini0288-velocitygraphlabel}
  \lean{PhyXMiniProblems.ProblemPhyXMini0288.VelocityGraphLabel}
  Named velocity levels and the time origin visible in the supplied graph
\end{definition}
\begin{proof}
  This is the declared model definition of Velocity Graph Label.
\end{proof}
\begin{definition}[Transverse Velocity Time Figure]
  \label{def:physics:phyx-mini-0288:phyxminiproblems-problemphyxmini0288-transversevelocitytimefigure}
  \lean{PhyXMiniProblems.ProblemPhyXMini0288.TransverseVelocityTimeFigure}
  \uses{def:physics:phyx-mini-0288:phyxminiproblems-problemphyxmini0288-signedlengthquantity, def:physics:phyx-mini-0288:phyxminiproblems-problemphyxmini0288-speedmagnitude, def:physics:phyx-mini-0288:phyxminiproblems-problemphyxmini0288-signedvelocityquantity, def:physics:phyx-mini-0288:phyxminiproblems-problemphyxmini0288-graphaxissymbol, def:physics:phyx-mini-0288:phyxminiproblems-problemphyxmini0288-velocitygraphlabel}
  The transverse-velocity-versus-time figure and its dimensionful labels
\end{definition}
\begin{proof}
  This is the declared model definition of Transverse Velocity Time Figure.
\end{proof}
\begin{definition}[Sinusoidal String Wave Setup]
  \label{def:physics:phyx-mini-0288:phyxminiproblems-problemphyxmini0288-sinusoidalstringwavesetup}
  \lean{PhyXMiniProblems.ProblemPhyXMini0288.SinusoidalStringWaveSetup}
  \uses{def:physics:phyx-mini-0288:phyxminiproblems-problemphyxmini0288-lengthmagnitude, def:physics:phyx-mini-0288:phyxminiproblems-problemphyxmini0288-signedlengthquantity, def:physics:phyx-mini-0288:phyxminiproblems-problemphyxmini0288-wavenumbermagnitude, def:physics:phyx-mini-0288:phyxminiproblems-problemphyxmini0288-angularfrequencymagnitude, def:physics:phyx-mini-0288:phyxminiproblems-problemphyxmini0288-speedmagnitude, def:physics:phyx-mini-0288:phyxminiproblems-problemphyxmini0288-signedvelocityquantity, def:physics:phyx-mini-0288:phyxminiproblems-problemphyxmini0288-signedaccelerationquantity, def:physics:phyx-mini-0288:phyxminiproblems-problemphyxmini0288-transversevelocitytimefigure}
  The physical wave setup. Time arguments are real coordinates measured in seconds. The position argument and all response values remain dimensionful. The phase is an angle modulo 2*pi; Real.Angle.toReal supplies its canonical representative in (-pi, pi]
\end{definition}
\begin{proof}
  This is the declared model definition of Sinusoidal String Wave Setup.
\end{proof}
\begin{definition}[Matches Transverse Velocity Graph]
  \label{def:physics:phyx-mini-0288:phyxminiproblems-problemphyxmini0288-matchestransversevelocitygraph}
  \lean{PhyXMiniProblems.ProblemPhyXMini0288.MatchesTransverseVelocityGraph}
  \uses{def:physics:phyx-mini-0288:phyxminiproblems-problemphyxmini0288-lengthinmeters, def:physics:phyx-mini-0288:phyxminiproblems-problemphyxmini0288-speedinmeterspersecond, def:physics:phyx-mini-0288:phyxminiproblems-problemphyxmini0288-velocityinmeterspersecond, def:physics:phyx-mini-0288:phyxminiproblems-problemphyxmini0288-accelerationinmeterspersecondsquared, def:physics:phyx-mini-0288:phyxminiproblems-problemphyxmini0288-sinusoidalstringwavesetup}
  Data read directly from the problem statement and primary image. The evenly spaced horizontal grid makes the unlabeled extrema +5 m/s and -5 m/s: the u\_s = 4 m/s line is one grid interval below the positive peak. At the vertical time axis the curve is at -u\_s and has positive slope. No field specifies the phase or selects an answer choice
\end{definition}
\begin{proof}
  This is the declared model definition of Matches Transverse Velocity Graph.
\end{proof}
\begin{definition}[Satisfies Sinusoidal Traveling Wave]
  \label{def:physics:phyx-mini-0288:phyxminiproblems-problemphyxmini0288-satisfiessinusoidaltravelingwave}
  \lean{PhyXMiniProblems.ProblemPhyXMini0288.SatisfiesSinusoidalTravelingWave}
  \uses{def:physics:phyx-mini-0288:phyxminiproblems-problemphyxmini0288-signedlengthquantity, def:physics:phyx-mini-0288:phyxminiproblems-problemphyxmini0288-lengthmagnitudereadout, def:physics:phyx-mini-0288:phyxminiproblems-problemphyxmini0288-signedlengthreadout, def:physics:phyx-mini-0288:phyxminiproblems-problemphyxmini0288-speedreadout, def:physics:phyx-mini-0288:phyxminiproblems-problemphyxmini0288-signedvelocityreadout, def:physics:phyx-mini-0288:phyxminiproblems-problemphyxmini0288-signedaccelerationreadout, def:physics:phyx-mini-0288:phyxminiproblems-problemphyxmini0288-lengthinmeters, def:physics:phyx-mini-0288:phyxminiproblems-problemphyxmini0288-wavenumberinradianspermeter, def:physics:phyx-mini-0288:phyxminiproblems-problemphyxmini0288-angularfrequencyinradianspersecond, def:physics:phyx-mini-0288:phyxminiproblems-problemphyxmini0288-sinusoidalstringwavesetup}
  The governing laws for the stated traveling-wave convention y(x,t) = y\_m sin (k*x - omega*t + phi). The derivative fields express transverse velocity as dy/dt and transverse acceleration as du/dt. The last two fields give the physical meaning of the maximum transverse speed. None of these laws states the requested numerical phase
\end{definition}
\begin{proof}
  This is the declared model definition of Satisfies Sinusoidal Traveling Wave.
\end{proof}
\begin{lemma}[transverse velocity normal form]
  \label{lem:physics:phyx-mini-0288:phyxminiproblems-problemphyxmini0288-transverse-velocity-normal-form}
  \lean{PhyXMiniProblems.ProblemPhyXMini0288.transverse_velocity_normal_form}
  \uses{def:physics:phyx-mini-0288:phyxminiproblems-problemphyxmini0288-signedlengthquantity, def:physics:phyx-mini-0288:phyxminiproblems-problemphyxmini0288-lengthmagnitudereadout, def:physics:phyx-mini-0288:phyxminiproblems-problemphyxmini0288-signedvelocityreadout, def:physics:phyx-mini-0288:phyxminiproblems-problemphyxmini0288-lengthinmeters, def:physics:phyx-mini-0288:phyxminiproblems-problemphyxmini0288-wavenumberinradianspermeter, def:physics:phyx-mini-0288:phyxminiproblems-problemphyxmini0288-angularfrequencyinradianspersecond, def:physics:phyx-mini-0288:phyxminiproblems-problemphyxmini0288-sinusoidalstringwavesetup, def:physics:phyx-mini-0288:phyxminiproblems-problemphyxmini0288-satisfiessinusoidaltravelingwave}
  Differentiating the source's displacement normal form gives u(x,t) = -omega*y\_m*cos(k*x - omega*t + phi) in coherent units. This is a derived law, not a premise field
\end{lemma}
\begin{proof}
  The conclusion follows from the governing relations and figure readouts named in the statement of transverse velocity normal form.
\end{proof}
\begin{lemma}[maximum transverse speed eq angular Frequency mul amplitude]
  \label{lem:physics:phyx-mini-0288:phyxminiproblems-problemphyxmini0288-maximum-transverse-speed-eq-angularfrequency-mul-amplitude}
  \lean{PhyXMiniProblems.ProblemPhyXMini0288.maximum_transverse_speed_eq_angularFrequency_mul_amplitude}
  \uses{def:physics:phyx-mini-0288:phyxminiproblems-problemphyxmini0288-lengthmagnitudereadout, def:physics:phyx-mini-0288:phyxminiproblems-problemphyxmini0288-speedreadout, def:physics:phyx-mini-0288:phyxminiproblems-problemphyxmini0288-angularfrequencyinradianspersecond, def:physics:phyx-mini-0288:phyxminiproblems-problemphyxmini0288-sinusoidalstringwavesetup, def:physics:phyx-mini-0288:phyxminiproblems-problemphyxmini0288-satisfiessinusoidaltravelingwave}
  The velocity amplitude of the sinusoidal wave is omega*y\_m
\end{lemma}
\begin{proof}
  The conclusion follows from the governing relations and figure readouts named in the statement of maximum transverse speed eq angular Frequency mul amplitude.
\end{proof}
\begin{lemma}[phase from transverse velocity graph]
  \label{lem:physics:phyx-mini-0288:phyxminiproblems-problemphyxmini0288-phase-from-transverse-velocity-graph}
  \lean{PhyXMiniProblems.ProblemPhyXMini0288.phase_from_transverse_velocity_graph}
  \uses{def:physics:phyx-mini-0288:phyxminiproblems-problemphyxmini0288-sinusoidalstringwavesetup, def:physics:phyx-mini-0288:phyxminiproblems-problemphyxmini0288-matchestransversevelocitygraph, def:physics:phyx-mini-0288:phyxminiproblems-problemphyxmini0288-satisfiessinusoidaltravelingwave}
  At x = 0, t = 0, the graph gives u/u\_max = -4/5. The derived velocity law therefore gives cos(phi) = 4/5. The graph is rising there, so its acceleration is positive; differentiating once more puts the canonical phase in quadrant IV. Thus its exact principal representative is -arccos (4/5)
\end{lemma}
\begin{proof}
  The conclusion follows from the governing relations and figure readouts named in the statement of phase from transverse velocity graph.
\end{proof}
\begin{definition}[Answer Choice]
  \label{def:physics:phyx-mini-0288:phyxminiproblems-problemphyxmini0288-answerchoice}
  \lean{PhyXMiniProblems.ProblemPhyXMini0288.AnswerChoice}
  The four phase readouts printed as answer choices, in radians
\end{definition}
\begin{proof}
  This is the declared model definition of Answer Choice.
\end{proof}
\begin{definition}[displayed Phase Radians]
  \label{def:physics:phyx-mini-0288:phyxminiproblems-problemphyxmini0288-displayedphaseradians}
  \lean{PhyXMiniProblems.ProblemPhyXMini0288.displayedPhaseRadians}
  \uses{def:physics:phyx-mini-0288:phyxminiproblems-problemphyxmini0288-answerchoice}
  Numerical radian readout printed beside each answer label
\end{definition}
\begin{proof}
  This is the declared model definition of displayed Phase Radians.
\end{proof}
\begin{definition}[recorded Dataset Answer]
  \label{def:physics:phyx-mini-0288:phyxminiproblems-problemphyxmini0288-recordeddatasetanswer}
  \lean{PhyXMiniProblems.ProblemPhyXMini0288.recordedDatasetAnswer}
  \uses{def:physics:phyx-mini-0288:phyxminiproblems-problemphyxmini0288-answerchoice}
  The answer label recorded in the source dataset
\end{definition}
\begin{proof}
  This is the declared model definition of recorded Dataset Answer.
\end{proof}
\begin{definition}[Is Nearest Displayed Phase]
  \label{def:physics:phyx-mini-0288:phyxminiproblems-problemphyxmini0288-isnearestdisplayedphase}
  \lean{PhyXMiniProblems.ProblemPhyXMini0288.IsNearestDisplayedPhase}
  \uses{def:physics:phyx-mini-0288:phyxminiproblems-problemphyxmini0288-answerchoice, def:physics:phyx-mini-0288:phyxminiproblems-problemphyxmini0288-displayedphaseradians}
  A choice is at least as close as every displayed phase readout
\end{definition}
\begin{proof}
  This is the declared model definition of Is Nearest Displayed Phase.
\end{proof}
% --- Archon declaration topology end ---
% --- Archon physics formalization source end ---
